\section{Enfoque Metodológico de Desarrollo e Ingeniería de Software}
\label{sec:metodologia_desarrollo}
% =============================================================================

La construcción del prototipo de votación electrónica se enmarca en las directrices de la investigación en ciencias del diseño (\textit{Design Science Research}, DSR) de Vaishnavi y Kuechler (2004) y Hevner et al. (2004), articulando un ciclo iterativo orientado a la resolución de problemas reales de seguridad de la información. Dentro de este marco científico, la ingeniería de la capa de aplicación web se instrumenta mediante la metodología \textit{UML-based Web Engineering} (UWE) propuesta por Koch et al. (2008), manteniendo una estricta separación de preocupaciones entre el método de investigación, el diseño de la interfaz web y la verificación criptográfica formal.

\subsection{Articulación Metodológica DSR--UWE e Ingeniería Criptográfica}
\label{subsec:articulacion_dsr_uwe}

El desarrollo del artefacto se estructuró a lo largo de las cinco fases metodológicas del ciclo general de DSR (Vaishnavi y Kuechler, 2004), integrando en las fases de sugerencia, desarrollo y evaluación los cinco modelos operativos de UWE (requisitos, contenido, navegación, presentación y procesos):

\begin{enumerate}
	\item \textbf{Conciencia del Problema (\textit{Awareness of Problem}):} Identificación de la vulnerabilidad crítica de los sistemas electorales contemporáneos frente a la computación cuántica (algoritmo de Shor) y a la recolección pasiva de boletas para descifrado futuro (\textit{Store-Now-Decrypt-Later}, HNDL), así como de los riesgos de suplantación y doble sufragio en interfaces web carentes de control formal de navegación, detallados en los Capítulos 1 y 2.
	\item \textbf{Sugerencia (\textit{Suggestion}):} Concepción de una arquitectura socio-técnica que acopla un protocolo híbrido post-cuántico —encapsulamiento asimétrico sobre retículos ML-KEM-768 (FIPS 203), firmas digitales ML-DSA-65 (FIPS 204), cifrado autenticado AES-256-GCM y recibos deterministas SHA3-256 (FIPS 202)— con un modelo de requisitos web UWE clasificados mediante estereotipos \texttt{«navigation»} y \texttt{«process»}.
	\item \textbf{Desarrollo (\textit{Development}):} Ingeniería sistemática del artefacto en dos frentes sincronizados: (a) el modelado UWE de la capa web (Modelo de Contenido acoplado al esquema relacional 3FN con disociación de identidad, Modelo de Navegación restringido por invariantes OCL, Modelo de Presentación abstracto y Modelo de Procesos web), y (b) la programación del núcleo criptográfico nativo en lenguaje C (\texttt{PQClean}), su enlace mediante wrappers Python FFI (\texttt{ctypes}) y la API REST en Flask.
	\item \textbf{Evaluación (\textit{Evaluation}):} Validación multidimensional mediante demostración matemática simbólica de cinco lemas en \textit{Tamarin Prover} 1.12.0 y \textit{Maude} 3.5.1, comprobación determinista de 200 vectores KAT oficiales del NIST, verificación de conformidad de los flujos de navegación y procesos UWE frente al prototipo (incluyendo la auditoría del inventario \texttt{PEND-01} a \texttt{PEND-22}) y pruebas experimentales de desempeño frente a la línea base clásica (RSA-2048/ECDSA P-256).
	\item \textbf{Conclusión (\textit{Conclusion}):} Consolidación de métricas de viabilidad operativa, síntesis de la contribución del modelado UWE en la trazabilidad de la seguridad y contraste formal de la hipótesis de investigación (\textbf{H1}).
\end{enumerate}

La Tabla \ref{tab:mapeo_dsr_uwe_cripto} formaliza la correspondencia estructural entre las fases del ciclo DSR, las actividades de modelado de la metodología UWE y los componentes de ingeniería criptográfica y verificación formal.

\begin{table}[!htbp]
	\centering
	\renewcommand{\arraystretch}{1.28}
	\footnotesize
	\caption{Mapeo metodológico entre las fases de DSR (Vaishnavi y Kuechler, 2004), el modelado UWE y la ingeniería criptográfica.}
	\label{tab:mapeo_dsr_uwe_cripto}
	\begin{tabularx}{\textwidth}{|>{\bfseries\raggedright\arraybackslash}p{2.6cm}|>{\raggedright\arraybackslash}p{4.3cm}|>{\raggedright\arraybackslash}p{4.3cm}|>{\raggedright\arraybackslash}X|}
		\hline
		\textbf{Fase DSR} & \textbf{Actividades en Ingeniería Web (UWE)} & \textbf{Actividades en Criptografía y Verificación Formal} & \textbf{Salidas y Artefactos Verificables} \\
		\hline
		1. Conciencia del Problema & Identificación de actores web ($V$, $AR$, $SV$, Auditor) y riesgos de navegación repetida o fuga de estado. & Análisis de amenaza cuántica (Shor/Grover) y formulación del modelo adversarial Dolev--Yao. & Definición de requerimientos RF1--RF4, RNF1--RNF2 y matriz de visibilidad de adversarios. \\
		\hline
		2. Sugerencia & Modelado de Requisitos Web UWE (\texttt{«navigation»} y \texttt{«process»}) y matriz de trazabilidad funcional. & Diseño del esquema híbrido KEM-DEM (ML-KEM-768 + AES-256-GCM con AAD) y firmas ML-DSA-65 con separación de dominio. & Diagrama de casos de uso UWE (Sección \ref{sec:casos_uso_actividades}) y diseño de ceremonia \textit{air-gapped} de $dk_{\text{AC}}$. \\
		\hline
		3. Desarrollo & Elaboración de los modelos UWE de Contenido, Navegación (con reglas OCL), Presentación y Procesos web. & Compilación de \texttt{libpqc\_core.dylib} en C, wrappers FFI, especificación MSR \texttt{protocolo\_v3.spthy} y línea base RSA/ECDSA. & Modelos UWE (Secciones \ref{sec:diagrama_clases_detallado}--\ref{sec:modelo_procesos_uwe}), esquema 3FN y prototipo funcional. \\
		\hline
		4. Evaluación & Verificación de conformidad funcional del prototipo respecto a los modelos UWE y auditoría de \texttt{PEND-15}. & Prueba automatizada de lemas en \textit{Tamarin Prover}, validación de vectores KAT NIST y benchmark de latencia. & Traza \texttt{salida\_tamarin.txt}, Tabla de Conformidad (\ref{tab:conformidad_prototipo_diseno}) y contraste estadístico de H1 (Cap.\ 4). \\
		\hline
		5. Conclusión & Síntesis del aporte de la ingeniería web dirigida por modelos en sistemas electorales críticos. & Delimitación de garantías probadas frente a límites simbólicos (\textit{diff-equivalence}, coacción) y hoja de ruta Wasm. & Conclusiones y recomendaciones técnicas para despliegue en navegador vía WebAssembly (Cap.\ 5). \\
		\hline
	\end{tabularx}
	\par\vspace{3pt}{\footnotesize\textit{Fuente: Elaboración propia basada en Vaishnavi y Kuechler (2004) y Koch et al. (2008).}}
\end{table}

\subsection{Requerimientos de Ingeniería del Prototipo}
\label{subsec:requerimientos_ingenieria}

El diseño del artefacto atiende a requerimientos funcionales (RF) y no funcionales (RNF) claramente acotados que vinculan la capa de interfaz web con el núcleo criptográfico:
\begin{itemize}
	\item \textbf{RF1 - Acreditación y Sufragio Único:} Todo votante debe contar con una credencial digital firmada por la Autoridad de Registro (AR) con contexto \texttt{ELEC\_CRED\_AR\_V1} antes de acceder a la cabina web de votación, garantizando que ninguna identidad emita más de un sufragio válido.
	\item \textbf{RF2 - Secreto del Voto y Disociación:} El contenido de la boleta debe permanecer cifrado en todo momento durante el transporte y custodia en el ánfora; la clave privada de desencapsulamiento ($dk_{\text{AC}}$) reside de forma exclusiva en la Autoridad de Conteo (AC) fuera de línea, y el almacenamiento relacional desacopla toda referencia entre el padrón y el voto cifrado.
	\item \textbf{RF3 - Integridad, Autenticidad y No Repudio Comicial:} Toda boleta transmitida desde el proceso web de emisión debe estar firmada digitalmente mediante ML-DSA-65 por la clave privada del elector ($sk_V$) vinculada a su credencial, impidiendo la alteración del voto en tránsito.
	\item \textbf{RF4 - Auditabilidad Individual y Recibo Inmutable:} Tras el depósito en el ánfora, la interfaz debe presentar y permitir descargar un comprobante firmado por el Servidor de Votación ($sk_{\text{SV}}$) que contenga el resumen determinista SHA3-256 de la boleta, verificable públicamente en el portal de auditoría sin revelar el sentido del voto.
	\item \textbf{RNF1 - Eficiencia y Rendimiento Operativo:} El tiempo de cifrado, firma y procesamiento por sufragio debe mantener una sobrecarga de latencia que no exceda en más del 30\% a la del protocolo clásico equivalente (RSA-2048/ECDSA P-256), garantizando una respuesta interactiva fluida en la cabina web.
	\item \textbf{RNF2 - Modularidad, Navegación Irreversible y Aislamiento de Memoria:} La capa de navegación web debe bloquear mediante restricciones formales toda transición de retorno a la cabina una vez emitido el voto, mientras que el núcleo nativo en C debe operar con buffers estáticos preasignados para el enlace FFI.
\end{itemize}

% =============================================================================
